\begin{proposition}[局域标架变换的联络律]
若 $\Phi'=D\Phi$、$D_\mu=\partial_\mu+\omega_\mu$，要求 $D'_\mu\Phi'=D(D_\mu\Phi)$，则\[\omega'_\mu=D\omega_\mu D^{-1}-(\partial_\mu D)D^{-1}.\]
\end{proposition}

\begin{proof}
把 $\Phi'=D\Phi$ 代入协变条件\eqnrefs{eq:ms-curved-D-cov} 两边，比较即可反解 $\omega'_\mu$。记 $\omega_\mu\equiv-\frac{i}{2}\omega_{\mu ab}\Sigma^{ab}$ 为 Lie 代数-取值 联络。左边由 Leibniz 法则展开
\begin{equation*}
D'_\mu\Phi'=(\partial_\mu D)\Phi+D\,\partial_\mu\Phi+\omega'_\mu\,D\Phi,
\end{equation*}
其中 $(\partial_\mu D)\Phi$ 是 $\partial_\mu$ 作用于乘积 $D(\Lambda(x))\Phi(x)$ 时对变换矩阵 $D$ 本身求导所得的额外项，正是 联络 要吸收的部分。右边由 $D$ 的线性性为
\begin{equation*}
D\,D_\mu\Phi=D\,\partial_\mu\Phi+D\,\omega_\mu\,\Phi.
\end{equation*}
两边相消 $D\,\partial_\mu\Phi$，得
\begin{equation*}
(\partial_\mu D)\Phi+\omega'_\mu\,D\Phi=D\,\omega_\mu\,\Phi.
\end{equation*}
把 $\omega'_\mu$ 项留在左边，右乘 $D^{-1}$ 并利用 $\Phi$ 任意性，即得
\begin{equation}\label{eq:ms-curved-omega-transform}
\boxed{
\omega'_\mu=D\,\omega_\mu\,D^{-1}-(\partial_\mu D)\,D^{-1}.
}
\end{equation}
写成 Lorentz 矩阵形式即 $\omega'_\mu=\Lambda\,\omega_\mu\Lambda^{-1}-(\partial_\mu\Lambda)\Lambda^{-1}$：第一项是伴随共轭，第二项 $(\partial_\mu\Lambda)\Lambda^{-1}$ 是非齐次项，补偿局域 标架 的变化。

这与电磁规范势 $A'_\mu=\Lambda A_\mu\Lambda^{-1}-(\partial_\mu\Lambda)\Lambda^{-1}$ 同构——自旋联络是局域 Lorentz 连接；对旋量须提升到 $\operatorname{Spin}(1,3)$，并由所选表示作用。

由此，自旋联络 与 Yang--Mills 势在规范群、变换、场强（曲率 $R_{\mu\nu}$ 对应 $F_{\mu\nu}$）各层平行；差异在于引力另有标架 $e_\mu^a$ 这一 soldering 形式，且 $\omega$ 由无挠条件 $T^a=de^a+\omega^a{}_b\wedge e^b=0$ 唯一确定为 Levi--Civita 自旋联络。沿闭合曲线的 平行运输 $W(C)=\mathcal P\exp(-\oint_C\omega_\mu\dd x^\mu)$ 给出 Lorentz 和乐，曲率即其无穷小生成元。
\end{proof}

\begin{proposition}[固定背景上的协变 Euler--Lagrange 方程]
对固定度量及固定丛联络，$S=\int\sqrt{|g|}\mathcal L(\Phi,D\Phi)\dd^dx$ 的紧支驻定场满足\[\frac{\partial\mathcal L}{\partial\Phi^A}-D_\mu\Pi_A^\mu=0.\]$D_\mu\Pi_A^\mu$ 包含时空指标的联络与内部对偶表示的联络。
\end{proposition}

\begin{proof}
与平直情形\eqnrefs{eq:ms-poincare-EL} 平行，只是用协变分部积分代替普通分部积分。场变分 $\Phi^A\to\Phi^A+\delta\Phi^A$ 给出
\begin{equation*}
\delta S=\int \dd^dx\sqrt{|g|}\left[\frac{\partial\mathcal{L}}{\partial\Phi^A}\delta\Phi^A+\frac{\partial\mathcal{L}}{\partial(D_\mu\Phi^A)}D_\mu\delta\Phi^A\right].
\end{equation*}
定义共轭动量 $\Pi_A^\mu\equiv\partial\mathcal{L}/\partial(D_\mu\Phi^A)$。协变分部积分用恒等式 $\int\dd^dx\sqrt{|g|}\,\nabla_\mu V^\mu=\int\dd^dx\,\partial_\mu(\sqrt{|g|}V^\mu)$，右端为全散度、积分后化为边界项；$\delta\Phi^A$ 与 $\Phi^A$ 同表示，联络项随之正确传递。忽略边界即得
\begin{equation*}
\int\dd^dx\sqrt{|g|}\,\Pi_A^\mu D_\mu\delta\Phi^A
=-\int\dd^dx\sqrt{|g|}\,(D_\mu\Pi_A^\mu)\delta\Phi^A.
\end{equation*}
代回，
\begin{equation*}
\delta S=\int\dd^dx\sqrt{|g|}\left[\frac{\partial\mathcal{L}}{\partial\Phi^A}-D_\mu\Pi_A^\mu\right]\delta\Phi^A+\text{(boundary)}.
\end{equation*}
边界项因变分在边界消失而不贡献，由 $\delta\Phi^A$ 任意性，
\begin{equation}\label{eq:ms-curved-EL}
\boxed{
\frac{\partial\mathcal{L}}{\partial\Phi^A}-D_\mu\frac{\partial\mathcal{L}}{\partial(D_\mu\Phi^A)}=0.
}
\end{equation}
形式与平直 Euler--Lagrange 方程\eqnrefs{eq:ms-poincare-EL} 一致，仅 $\partial_\mu\to D_\mu$。
\end{proof}

\begin{proposition}[微分同胚不变性与协变守恒]
若物质作用量微分同胚不变，所有参与变换的物质场满足方程，则按 $T_{\mu\nu}=-2|g|^{-1/2}\delta S_m/\delta g^{\mu\nu}$ 定义的 Hilbert 张量满足 $\nabla_\mu T^{\mu\nu}=0$。有外部背景场时须加入相应力密度。
\end{proposition}

\begin{proof}
把无穷小 微分同胚 $x^\mu\to x^\mu+\xi^\mu(x)$ 引起的度规 Lie 导数 $\delta_\xi g_{\mu\nu}=\nabla_\mu\xi_\nu+\nabla_\nu\xi_\mu$ 代入 Hilbert 张量定义\eqnrefs{eq:ms-curved-T-hilbert}。由 $\delta_g S_m=+\frac12\int\dd^dx\sqrt{|g|}\,T^{\mu\nu}\delta g_{\mu\nu}$，
\begin{equation*}
\delta_\xi S_m=+\frac12\int\dd^dx\sqrt{|g|}\,T^{\mu\nu}(\nabla_\mu\xi_\nu+\nabla_\nu\xi_\mu).
\end{equation*}
Hilbert 张量自动对称（$\delta g^{\mu\nu}=\delta g^{\nu\mu}$），交换哑指标 $\mu\leftrightarrow\nu$ 后两项相等，故
\begin{equation*}
\delta_\xi S_m=+\int\dd^dx\sqrt{|g|}\,T^{\mu\nu}\nabla_\mu\xi_\nu.
\end{equation*}
协变分部积分。由 Leibniz 法则 $T^{\mu\nu}\nabla_\mu\xi_\nu=\nabla_\mu(T^{\mu\nu}\xi_\nu)-(\nabla_\mu T^{\mu\nu})\xi_\nu$，第一项是全协变散度，用恒等式 $\int\dd^dx\sqrt{|g|}\,\nabla_\mu V^\mu=\int\dd^dx\,\partial_\mu(\sqrt{|g|}V^\mu)$ 化为边界项；$\xi^\nu$ 在边界消失，故
\begin{equation*}
\delta_\xi S_m=-\int\dd^dx\sqrt{|g|}\,(\nabla_\mu T^{\mu\nu})\xi_\nu.
\end{equation*}
在壳 时 微分同胚 不变性 给出 $\delta_\xi S_m=0$。$\xi^\nu(x)$ 任意，故被积函数系数逐点为零：
\begin{equation}\label{eq:ms-curved-T-conservation}
\boxed{\nabla_\mu T^{\mu\nu}=0.}
\end{equation}
这就是协变守恒律，与平直 $\partial_\mu T^{\mu\nu}=0$ 相比仅以协变导数取代普通导数。

因 $\nabla_\mu$ 含联络项，此式一般不能直接积分为守恒荷；需 Killing 矢量 $K^\mu$（$\nabla_{(\mu}K_{\nu)}=0$）使 $J_K^\mu=T^{\mu\nu}K_\nu$ 散度为零，才得守恒荷 $Q_K=\int_\Sigma T^{\mu\nu}K_\nu\,\dd\Sigma_\mu$：Minkowski 的 10 个 Killing 矢量给出能量-动量与角动量守恒。

引力没有一般坐标不变的局域能量密度；全局或准局域能量须结合额外结构定义，例如渐近平直空间的 ADM 质量、零无穷远的 Bondi 质量或特定对称背景下的 Komar 荷。
\end{proof}

\begin{proposition}[有质量 Proca 场的横向约束]
在 $(-,+,+,+)$ 约定下，$\partial_\mu F^{\mu\nu}-m^2A^\nu=0$、$m\ne0$ 推出 $\partial_\mu A^\mu=0$。对 $p^2=-m^2$ 的平面波，极化属于三维空间 $p^\perp$；这是运动方程推论而非规范条件。
\end{proposition}

\begin{proof}
对 Proca 运动方程 $\partial_\mu F^{\mu\nu}-m^2A^\nu=0$ 取 $\partial_\nu$：
\begin{equation*}
\partial_\nu\partial_\mu F^{\mu\nu}-m^2\partial_\nu A^\nu=0.
\end{equation*}
第一项恒为零：$\partial_\nu\partial_\mu$ 对 $\mu,\nu$ 对称（偏导数可交换），而 $F^{\mu\nu}$ 对 $\mu,\nu$ 反对称，对称张量与反对称张量缩并为零：
\begin{equation*}
\partial_\nu\partial_\mu F^{\mu\nu}
=\tfrac12(\partial_\nu\partial_\mu+\partial_\mu\partial_\nu)F^{\mu\nu}=0,
\end{equation*}
末步因交换哑指标 $\mu\leftrightarrow\nu$ 后 $F^{\nu\mu}=-F^{\mu\nu}$，两项相消。于是 $-m^2\partial_\nu A^\nu=0$；由 $m\neq0$ 得 Lorenz 条件
\begin{equation*}
\partial_\mu A^\mu=0,
\end{equation*}
它是运动方程的推论而非额外规范选择。

自由度计数：$A^\mu$ 的 4 个 Lorentz 分量被一个标量约束 $\partial_\mu A^\mu=0$ 降为 3 个独立极化 $\lambda=-1,0,+1$，即 有质量 自旋-1 的 小群 $SO(3)$ 的 $2s+1=3$ 个态；$\lambda=\pm1$ 为横向、$\lambda=0$ 为纵向。

当 $m=0$ 时本推导失效（需 $m\neq0$），此时 规范 冗余 $A_\mu\to A_\mu+\partial_\mu\alpha$ 额外消除纵向自由度，只剩 $\lambda=\pm1$；$m\to0$ 的纵向模不能仅凭规范语言删去：自由 Stueckelberg 理论的平滑极限一般包含 Maxwell 场和解耦标量。

电弱破缺中则由规范场与 Goldstone 模耦合产生有质量矢量的纵向极化。

Stueckelberg 机制把 Proca 改写为形式规范不变的 $\mathcal{L}_{\mathrm{St}}=-\tfrac14F_{\mu\nu}F^{\mu\nu}-\tfrac12m^2(A_\mu-\partial_\mu\phi)^2$，幺正规范 $\phi=0$ 恢复 Proca，$\phi$ 即被吃掉的 Goldstone 模；纵向极化 $\epsilon^{(0)}\sim p_\mu/m$ 在高能主导，正是等效定理的运动学根源。
\end{proof}
